  \section{Introduction}
  \label{sec:introduction}

  \subsection{Digital twins and language models}
  \label{sec:background}

  A digital twin (DT) is a software system that mirrors a running physical system, with data
  moving automatically in both directions between the two
  \cite{fuller2020digitaltwin,kritzinger2018digitaltwin}. It is characterised by the services it
  offers on top of that connection, such as monitoring, decision support and reconfiguration, and
  it is not a control system: the physical system must keep working if the DT is switched off
  \cite{gomes2025incubator}. This project applies the same principle one level higher up, adding a
  layer on top of a DT that must keep working without that layer.

  A large language model (LLM) is a neural network trained on large quantities of text that can
  carry out a task described in ordinary language without being trained for it
  \cite{brown2020fewshot}. Its output can be fluent and still wrong, stating false quantities with
  confidence, a failure known as hallucination \cite{ji2023hallucination}. It can also call
  external code and continue with the returned value \cite{schick2023toolformer}, which shifts its
  responsibility from knowing a quantity to knowing which computation produces it.

  The two are complementary: a DT holds a calibrated model of a process, while an LLM is easy to
  address but has no reliable account of how any particular process behaves. \Cref{sec:related}
  reviews the work that builds on this pairing.


  \subsection{The incubator DT}
  \label{sec:incubator-dt}

  This project builds on the incubator DT of Gomes et al. \cite{gomes2025incubator}: an insulated
  box with a heating element, a fan and temperature sensors, and a controller that switches the
  heater to hold the air at a setpoint. The box is used to ferment tempeh, a food made by growing a
  mould on cooked soybeans over roughly a day, whose success depends on the temperature staying
  within a fairly narrow range \cite{steinkraus1960tempeh,tempehinfo_temperature}. The DT consists
  of thermal models of the box, calibrated against its own measurements and kept supplied with live
  data, and offers three groups of services on top of them \cite{gomes2025incubator}:

  \begin{itemize}
    \item \textbf{Visualisation and monitoring.} It presents the current and past state and
      compares what the models predict with what the sensors report, estimating unmeasured
      quantities from measured ones. A growing discrepancy signals that something has changed.
    \item \textbf{Decision support.} It simulates scenarios before they happen, known as what-if
      analysis, and searches across simulated configurations for the best trade-off. Gomes et al.
      give two examples: heating the box above its setpoint before a power supply is replaced, with
      simulation choosing how long to heat so that the batch neither overheats nor cools too far
      during the outage, and mapping how controller settings trade tight temperature control
      against heater wear.
    \item \textbf{Reconfiguration.} When the models no longer match the hardware, it recalibrates
      them against fresh data and re-tunes the controller.
  \end{itemize}

  None of these services is modified by this work.

  \subsection{The gap this project addresses}
  \label{sec:gap}

  Consider an operator who wants to open the lid of a running batch for five minutes, and wants to
  know beforehand whether the batch can take it. The DT can answer this in principle. It has
  calibrated thermal models, a representation of the disturbance an open lid causes, and a
  decision-support service meant for this kind of before-the-fact analysis, the power supply
  example is structurally the same problem \cite{gomes2025incubator}. The gap is not a missing
  capability but the distance between that capability and the person who needs it:

  \begin{itemize}
    \item \textbf{Using the service means building the experiment.} A what-if analysis is a
      simulation someone configures, choosing the model, the parameters that represent the
      disturbance, the horizon and what to extract. An operator who knows fermentation but not
      thermal parameters has no route to it.
    \item \textbf{The service stops at numbers.} It returns a predicted temperature over time,
      while the operator needs a judgement about the batch. That requires knowing what the culture
      tolerates, for how long and at which stage, and the DT does not hold that knowledge.
    \item \textbf{Every question has to be anticipated.} The published examples exist because
      someone foresaw those situations and wrote code for them, so what the DT can be asked is
      fixed at development time.
  \end{itemize}

  This project investigates a natural-language advisory layer on top of these services, which turns
  the operator's question into a correctly configured call and the result back into advice about
  the batch, without changing how the DT computes.

  \subsection{The LLM advisory layer architecture}
  \label{sec:architecture}

  The advisory is organised as three ordered layers (\Cref{fig:architecture}). A layer may call the
  one below it and never the one above, so a quantity reaches the operator only after something
  below has computed it.

  \begin{figure}[!ht]
    \centering
    \setstretch{1}
    \definecolor{archteal}{HTML}{1F7A73}
    \definecolor{archblue}{HTML}{2F6DB5}
    \definecolor{archsand}{HTML}{F3E7CF}
    \resizebox{\textwidth}{!}{%
    \begin{tikzpicture}[
        box/.style={draw=black!70, fill=white, rounded corners=2pt, align=center,
                    font=\small, minimum height=1.2cm, inner sep=3pt},
        src/.style={box, minimum height=1.6cm},
        ours/.style={fill=archteal!10},
        dt/.style={fill=archsand},
        cont/.style={draw=black!45, rounded corners=4pt},
        flow/.style={-{Stealth[length=2.2mm]}, semithick, draw=archteal},
        bflow/.style={flow, draw=archblue, dashed},
        cflow/.style={-{Stealth[length=2.2mm]}, semithick, draw=black!60, densely dotted},
        lab/.style={font=\scriptsize, text=archteal, align=center},
        blab/.style={font=\scriptsize, text=archblue, align=center},
        band/.style={draw=black!20, rounded corners=3pt},
        bandname/.style={rotate=90, font=\footnotesize\scshape, text=black!65},
        note/.style={font=\scriptsize\itshape, text=black!55, align=center},
        swatch/.style={draw=black!45, minimum width=4.5mm, minimum height=3mm, inner sep=0pt},
        legtext/.style={font=\scriptsize, inner sep=1pt}]

      % Bands
      \draw[band] (-0.9,8.6) rectangle (15.1,13.1);
      \draw[band] (-0.9,3.2) rectangle (15.1,8.0);
      \draw[band] (-0.9,0.0) rectangle (15.1,2.6);
      \node[bandname] at (-0.5,10.85) {Language layer};
      \node[bandname] at (-0.5,5.6) {Tool layer};
      \node[bandname] at (-0.5,1.3) {Sources};
      \node[note] at (13.8,10.85) {orchestrates\\and explains,\\predicts and\\judges nothing};
      \node[note] at (2.2,6.9) {computes every prediction\\and verdict};

      % Language layer: the skills and the harness that runs the model
      \draw[cont, fill=archteal!10] (0.0,8.9) rectangle (3.4,12.8);
      \node[font=\footnotesize\bfseries] at (1.7,12.45) {Skills};
      \node[note] at (1.7,9.28) {each a description\\and instructions};
      \draw[cont, fill=black!9] (3.8,8.9) rectangle (12.5,12.8);
      \node[anchor=north east, font=\footnotesize\bfseries] at (12.4,12.75) {Claude Code harness};

      % Nodes
      \node[box, minimum width=3cm, rounded corners=8pt, font=\normalsize] (op) at (5.6,13.9) {Operator};
      \node[box, minimum width=2.9cm, minimum height=0.9cm] (lidsk) at (1.7,11.5) {Lid-opening\\skill};
      \node[box, minimum width=2.9cm, minimum height=0.9cm] (hbsk)  at (1.7,10.3) {Heater-boost\\skill};
      \node[box, minimum width=2.8cm, minimum height=2.1cm] (llm) at (5.6,10.85)
            {LLM\\{\scriptsize\itshape reads the question,}\\{\scriptsize\itshape chooses each command,}\\{\scriptsize\itshape writes the answer}};
      \node[box, minimum width=3.4cm, minimum height=2.1cm] (tex) at (10.4,10.85)
            {Tool execution\\{\scriptsize\itshape runs each command}\\{\scriptsize\itshape and returns its output}};
      \node[box, ours, minimum width=2.3cm] (st)   at (1.4,4.6)  {State\\access};
      \node[box, ours, minimum width=2.3cm] (sim)  at (5.2,4.6)  {Simulation};
      \node[box, ours, minimum width=2.3cm] (risk) at (9.0,4.6)  {Risk\\assessment};
      \node[box, ours, minimum width=2.6cm] (bo)   at (13.0,4.6) {Boost\\optimisation};
      \node[src, dt, minimum width=3.2cm] (dts) at (1.7,1.3)  {DT state\\{\scriptsize\itshape as published}};
      \node[src, dt, minimum width=3.8cm] (dtm) at (6.4,1.3)  {DT models\\and controllers\\{\scriptsize\itshape used unmodified}};
      \node[src, ours, minimum width=5.0cm] (dk)  at (12.3,1.3) {Domain knowledge\\{\scriptsize\itshape sourced or estimated}\\{\scriptsize phases, bands, budgets, heater limits}};

      % Attachment points
      \coordinate (texA)  at ([xshift=-14mm]tex.south);
      \coordinate (texB)  at ([xshift=-6mm]tex.south);
      \coordinate (texC)  at ([xshift=6mm]tex.south);
      \coordinate (texD)  at ([xshift=14mm]tex.south);
      \coordinate (riskU) at ([xshift=8mm]risk.north);
      \coordinate (boU)   at ([xshift=-4mm]bo.north);
      \coordinate (boV)   at ([xshift=6mm]bo.north);
      \coordinate (simM)  at ([xshift=-5mm]sim.south);
      \coordinate (simD)  at ([xshift=4mm]sim.south);
      \coordinate (simC)  at ([xshift=9mm]sim.south);
      \coordinate (riskL) at ([xshift=-5mm]risk.south);
      \coordinate (riskD) at ([xshift=10mm]risk.south);
      \coordinate (boL)   at ([xshift=-4mm]bo.south);
      \coordinate (boD)   at ([xshift=8mm]bo.south);
      \coordinate (y285)  at (0,2.85);
      \coordinate (y355)  at (0,3.55);
      \coordinate (y66)   at (0,6.6);
      \coordinate (y70)   at (0,7.0);
      \coordinate (y76)   at (0,7.6);

      % Operator, skills and the model's loop through the harness
      \draw[flow] ([xshift=-4mm]op.south) -- node[lab, left, pos=0.55] {question} ([xshift=-4mm]llm.north);
      \draw[flow] ([xshift=4mm]llm.north) -- node[lab, right, pos=0.45] {answer} ([xshift=4mm]op.south);
      \draw[cflow] (3.4,10.85) -- (llm.west);
      \draw[flow] ([yshift=3mm]llm.east) to[bend left=35] node[lab, above] {command} ([yshift=3mm]tex.west);
      \draw[flow] ([yshift=-3mm]tex.west) to[bend left=35] node[lab, below] {JSON\\result} ([yshift=-3mm]llm.east);
      \node[text=archteal, font=\large] at ($(llm.east)!0.5!(tex.west)$) {$\circlearrowright$};

      % Lid-opening tool: the pipeline
      \draw[flow] (texA) -- (texA |- y66) -| node[lab, above, pos=0.25] {duration, start time} (sim.north);
      \draw[flow] (riskU) -- node[lab, right, pos=0.18] {verdict and\\prediction} (texB);
      \draw[flow] (st) -- node[lab, above] {conditions} (sim);
      \draw[flow] (sim) -- node[lab, above] {trajectory} (risk);

      % Heater-boost tool and its coded loop
      \draw[bflow] (texC) -- (texC |- y70) -| node[blab, below, pos=0.25] {duration} (boU);
      \draw[bflow] (boV) -- node[blab, right] {verdict\\and plan} (boV |- y76) -| (texD);
      \draw[bflow] (boL) -- (boL |- y355) -| (simD);
      \draw[bflow] (riskL |- y355) -- (riskL);
      \fill[archblue] (riskL |- y355) circle (1.2pt);
      \node[text=archblue] at (14.65,4.6) {$\circlearrowright$};

      % Sources
      \draw[flow] (st.south |- dts.north) -- (st.south);
      \draw[flow] (simM |- dtm.north) -- (simM);
      \draw[flow] (riskD |- dk.north) -- (riskD);
      \draw[flow] (riskD |- y285) -| (simC);
      \fill[archteal] (riskD |- y285) circle (1.2pt);
      \draw[bflow] (boD |- dk.north) -- (boD);

      % Legend
      \draw[flow] (-0.9,-0.6) -- ++(0.9,0) node[right, font=\scriptsize, text=black] {every request};
      \draw[bflow] (2.4,-0.6) -- ++(0.9,0) node[right, font=\scriptsize, text=black]
                  {heater-boost requests only, each candidate boost is simulated and assessed in turn};
      \node[swatch, fill=black!9] (sw1) at (-0.675,-1.2) {};
      \node[legtext, right=1mm of sw1] (lg1) {Claude Code};
      \node[swatch, fill=archteal!10, right=5mm of lg1] (sw2) {};
      \node[legtext, right=1mm of sw2] (lg2) {this project};
      \node[swatch, fill=archsand, right=5mm of lg2] (sw3) {};
      \node[legtext, right=1mm of sw3] (lg3) {existing DT};
      \draw[cflow] ([xshift=6mm]lg3.east) -- ++(0.9,0) node[right, font=\scriptsize, text=black] {loaded into the model's context};
    \end{tikzpicture}}
    \caption{The advisory layer in three ordered layers, shaded by who provides each part. The
      model works through the harness in a loop ($\circlearrowright$), one command and its result at
      a time, and the heater-boost tool runs a coded loop of its own.}
    \label{fig:architecture}
  \end{figure}

  The \textbf{language layer} is an LLM working in Claude Code, an agent harness, with
  one \emph{skill} per advisory. A skill is a document of instructions whose description tells the model
  when to load it, and whose body is loaded only when it is used \cite{anthropic2026skills}. The
  model recognises which advisory a question belongs to, converts what the operator said into the
  tool's parameters and issues a command, which the harness runs, returning the tool's JSON
  output. The model may issue further commands before it answers. It orchestrates the tools and
  explains their results, but predicts and judges nothing.

  The \textbf{tool layer} computes every prediction and verdict, as a pipeline. State access assembles the
  batch's current conditions, the simulation projects them forward under the disturbance asked
  about, and the risk assessment reduces the trajectory to a verdict and the numbers and thresholds that
  decided it. The verdict crosses back to the language layer together with the prediction, because
  turning a curve into a decision about the batch must not be left to prose. The second advisory, a heater boost, reuses the same pipeline as an optimisation: each
  candidate pre-heating strategy is scored by re-running the simulation and the risk assessment,
  in a loop coded in the tool.

  The \textbf{sources} differ in where their authority comes from. The live DT state makes an
  answer concern this batch rather than batches in general. The DT's own models and controllers
  are driven from outside and never altered, so numerical authority stays with the DT
  \cite{gomes2025incubator}. Domain knowledge supplies the fermentation thresholds.

  \subsection{Research question}
  \label{sec:rq}

  This project takes up the third point of the gap (\Cref{sec:gap}). The twin's tools answer
  forward questions: given an opening of a set length, what happens to the batch. The questions an
  operator asks next are often inverse ones, such as how long the lid can stay open or when to open
  it. No tool answers these directly, but a search over the forward tools can, and it can be
  provided in two ways. It can be coded as a new tool, as the heater-boost advisory does for its
  setpoint (\Cref{sec:boost}), which means foreseeing the question. Or it can be left to the
  LLM, which runs the existing tools itself when the question is asked. The research
  question asks which of the two a twin should use:

  \begin{quote}
    When a DT needs to answer a new kind of question, should the answer be coded as a
    new tool or delegated to an LLM over the existing tools?
    \begin{enumerate}[label=(\alph*)]
      \item If it is delegated, how close are the model's answers to how-long and when questions
        to a search coded over the same tools, and does it say which parts the tools did not
        compute?
      \item What does delegating cost and save compared with coding: accuracy at the simulation's
        resolution, tokens and time per answer, and the need to foresee the question?
    \end{enumerate}
  \end{quote}
